\scanpage{193}
% archive sheet number: 192

types \deleted{simples}\inserted{``élémentaires''}, de part et d'autre,
\[
\unsure{\mathbb G_m|G_m},\quad\mathbb Z/\ell\mathbb Z,\quad\mu_p\quad\unsure{\mathbb G_a|G_a},\quad\alpha_p,\quad\mathbb Z/p\mathbb Z
\]
\margin{$\mu_p,\alpha_p,\mathbb Z/p\mathbb Z$ \unsure{spécial à} car. $p>0$}
[Et, l'ordre où on les rencontre dans une suite de composition canonique \unsure{est} $=$, à très peu de chose près, sur un corps non alg. clos ni $=$ parfait]
% slanted
\margin{Si on se borne à l'unipotent fini, on le \unsure{trouve en prenant|trouve en faisant} \unsure{relativement|des extensions} en bas [le \uline{multiplicatif} \unsure{correspond aux} discrets séparables]. Or Cartier a \deleted{\unsure{exactement}} classifié les \deleted{hyperalgèbres}\inserted{\unsure{grp. alg. finis}} \uline{unipotents}}
\[
\mu_p,\quad\alpha_p,\quad\mathbb V,\quad\mathbb Z/p\mathbb Z,\quad\mathbb Z/\ell\mathbb Z,\quad\mathbb Z
\]
$\mathbb V=D(\mathbb G_a)$ est le \unsure{cinquième} groupe; en car $p>0$, on a envie de les mettre ensemble en faisant des \uline{extensions}, \deleted{\unsure{Prenons une}}\inserted{dans la} \deleted{\unsure{propre}} catégorie (présumée abélienne) des groupes (ind-affines) \deleted{[\unsure{formels}]} [ind-$S$-schémas affines, avec des structures dessus]
% slanted
\margin{\unsure{Regarder} en général (sans condition de commutativité), \deleted{\unsure{pro-catégorie}} des ``groupes \deleted{algébriques} proalgébriques'' (extension d'un groupe formel par un \unsure{pro-grpe algébrique})]}

Un tel groupe aura une algèbre affine, qui sera une \deleted{limite projective \unsure{stricte}}\inserted{somme directe \inserted{infinie}} de \unsure{produits} de \unsure{discrets} (\unsure{Bo---|Bour---}) \deleted{\unsure{par}} (c'est un \unsure{$\beta$|$f$} \unsure{venant} d'un dans une \unsure{limite}, \unsure{un peu} dans le cas non \unsure{proalgébrique})

\deleted{\unsure{Opération} de $F,V$},
% slanted
\margin{dans le cas $k$ parfait, extensions triviales, \unsure{du moins} \unsure{le sont pour} \illegible{2}}
\diagram{\[\begin{array}{c@{\qquad}c@{\qquad}c}
\text{Multiplicatif [\unsure{quelle compt. est. connexe})}
&\text{unipot\deleted{\unsure{ent}}}
&\text{[\unsure{sans $\ell$-torsion ou $p$-torsion}]}\\
&\vert\!\rule{34mm}{0.5pt}\!\vert
&\text{groupes discrets \unsure{ét.}}\\
\text{Tore}\;\backslash\;\text{grpe profini \deleted{\unsure{étale}}}\;\times\mu_p
&&\vert\!\rule{40mm}{0.4pt}\!\vert\\
\text{\hspace{15mm}d'ordre premier à $p$}
&\text{\deleted{\unsure{groupe}} admet un}
&\text{grpe \deleted{\unsure{de torsion}}\inserted{\unsure{ordinaire}}}\\
\text{\unsure{dual aux} groupes \deleted{\unsure{libres}}\inserted{\unsure{discrets}}}
&\text{plus grand sous grpe alg. et}
&\text{\deleted{\unsure{à $\pi$ opère}}}\\
\text{\unsure{sans torsion à}\;$\unsure{\pi}$}
&\text{un plus gd \deleted{\unsure{partie}}\inserted{\unsure{quotient}}}
&\text{\unsure{où} $\pi$ opère}\\
\text{$\vert$ \unsure{dual aux} groupes de torsion première à $p$}
&\text{infinitésimal (sur $k$)}
&\text{continûment}\\
\text{\unsure{où opère} $\pi$ $\vert$ \unsure{dual aux grpes}}&&\\
\text{\unsure{de $p$-torsion où opère} $\pi$}&&
\end{array}\]}
